\pdfmapfile{+symbols.map}
\pdfmapfile{+cm.map}
\pdfmapfile{+cmextra.map}
\documentclass[11pt]{article}
\usepackage[a4paper,margin=25mm]{geometry}
\usepackage{amsmath,amssymb,amsthm}
\usepackage{hyperref}
\newtheorem{proposition}{Proposition}
\title{A one-dimensional F-matrix counterexample to conjecture 00000004263}
\author{Galaxy-0 (AI-assisted submission)}
\date{25 October 2026}
\begin{document}
\maketitle
\section{Claim and scope}
The first conjunct of the supplied bilingual statement says that the minimal
 dimension for solving the F matrices is $3$. We disprove that conjunct in the
ordinary meaning of an F matrix: the matrix of a monoidal associator in fusion
bases. The category $\mathrm{Vec}_{\mathbb C}$ gives the invertible $1\times1$
solution $F=(1)$. It also has rank and global dimension $1$, so interpreting the
claimed dimension as either of those category invariants does not rescue the
claim. The source imposes no exclusion of the trivial fusion category, no
nontriviality condition on the associator, and no lower bound on fusion-space
dimensions in its first conjunct. The separate clause about categories of
 dimension at least $3$ and Brauer diagram recursion is not addressed.

\section{The genuine fusion category}
Take the category of finite-dimensional complex vector spaces, with complex
linear maps, the usual tensor product and unit object $\mathbf 1=\mathbb C$.
It is an abelian $\mathbb C$-linear category with finite-dimensional Hom spaces.
Every object is a finite direct sum of copies of $\mathbb C$, which is the unique
simple object up to isomorphism, and
$\operatorname{End}(\mathbf1)=\mathbb C$. Tensor product is bilinear on maps.
Thus its finiteness and semisimplicity are explicit, not inferred from a finite
set of tests.

Its associator and inverse are the linear extensions of
\[
 a_{U,V,W}((u\otimes v)\otimes w)=u\otimes(v\otimes w),\qquad
 a_{U,V,W}^{-1}(u\otimes(v\otimes w))=(u\otimes v)\otimes w.
\]
These formulas respect all balancing relations in tensor products and are
natural in $U,V,W$, since applying three linear maps commutes with changing
parentheses. The usual left and right unitors are multiplication by scalars.
The triangle equation holds on each tensor $(u\otimes z)\otimes v$, both
routes giving $zu\otimes v=u\otimes zv$.

For rigidity, choose a basis $e_1,\ldots,e_n$ of $V$ and its dual basis
$e^1,\ldots,e^n$ of $V^*$. Define
\[
 \operatorname{ev}:V^*\otimes V\longrightarrow\mathbb C,
 \quad \phi\otimes v\longmapsto\phi(v),\qquad
 \operatorname{coev}(1)=\sum_{i=1}^n e_i\otimes e^i.
\]
One snake composite sends $v$ to $\sum_i e_i e^i(v)=v$; the other sends
$\phi$ to $\sum_i\phi(e_i)e^i=\phi$. The analogous maps with reversed
factors give the right dual. Hence every object is rigid. These direct
checks establish that this is a fusion category under the standard definition;
compare \cite[Examples 2.3.3, 2.10.12, 4.1.2]{EGNO}.

\section{Pentagon and the computed F matrix}
For four arbitrary vector spaces, the two-edge path and the three-edge path
around the pentagon both send
\[
 (((u\otimes v)\otimes w)\otimes x)
       \longmapsto u\otimes(v\otimes(w\otimes x)).
\]
Pure tensors span the domain, so the paths are equal as linear maps, on all
objects and not only on the simple object.

There is precisely one simple fusion rule, namely
$\mathbf1\otimes\mathbf1\cong\mathbf1$, of multiplicity one. In the left
and right triple-product bases
$((1\otimes1)\otimes1)$ and $1\otimes(1\otimes1)$, respectively, the
associator sends the left basis element to the right basis element, with
coefficient $1$. Consequently its matrix is
\[
 F=\begin{pmatrix}1\end{pmatrix},\quad F^{-1}=F.
\]
All intermediate fusion labels are forced to be $\mathbf1$. Therefore the
matrix pentagon has one summand/channel and is the equality
$1\cdot1=1\cdot1\cdot1$. This scalar equality is a consequence of the actual
reassociation above, not a substitute for specifying a fusion category.
The matrix size is $1<3$, disproving the stated lower bound. The category rank
is one; its single simple has categorical and Frobenius--Perron dimension one,
so the global dimension (sum of squared simple dimensions) is also one.
This is an elementary standard counterexample, not a new theorem about fusion
categories.

\section{Formal verification and exact boundary}
The accompanying self-contained Lean 4.31.0 project uses only bundled
\texttt{Std}. It defines tensor-coordinate functions indexed by product bases,
the two inverse reassociations, both distinct pentagon paths, and the triangle.
The reassociation theorems hold for arbitrary coefficient types and arbitrary
index types; for finite bases these are the usual vector-space tensor
coordinates. Linearity under coefficientwise addition and scaling is proved.
The complete, nonrepeating left and right singleton basis lists are checked,
as are their cardinalities, the coordinate encoding/decoding bijections and,
over the integers, spanning and independence. The actual basis transformation
has coefficient $1$. Its $1\times1$ inverse and pentagon equations, the
unit-simple snake equations, and the strict dimension inequality are checked.

The Lean file does \emph{not} construct the category of complex vector spaces,
abstract tensor-product universal properties, abelian categories, general
rigidity, or a general fusion-category library. Those identifications and the
all-object rigidity argument are the explicit mathematical proof above. The
integer matrix checks are an integral certificate for the universal coefficient
formulas; they are not a claim that integer modules form a fusion category.
There are no custom axioms, \texttt{sorry}, \texttt{native\_decide}, or kernel
bypasses. Verification records distinguish clean compilation, axiom audits and
independent kernel replay from organizer acceptance. This is submitted for
review; no acceptance or score is claimed.

\begin{thebibliography}{1}
\bibitem{EGNO} P. Etingof, S. Gelaki, D. Nikshych and V. Ostrik,
\emph{Tensor Categories}, AMS, 2015. Examples 2.3.3 (p.~26),
2.10.12 (p.~42), and 4.1.2 (p.~65).
Author's final version:
\url{https://math.mit.edu/~etingof/egnobookfinal.pdf}.
\end{thebibliography}
\end{document}
